% ============================================================================
\section{Agent-specific safety problems}
% ============================================================================

\begin{frame}{What Makes Agents Harder Than Models}
  \begin{itemize}\itemsep0.5em
    \item[\textcolor{KAred}{$\blacktriangleright$}] The same model is safer as a chatbot than as an agent. What changes is the privileges and the loop wrapped around it, not the weights.
    \item[\textcolor{KAred}{$\blacktriangleright$}] Four properties turn ordinary model mistakes into real-world harm:
      \begin{itemize}\itemsep0.3em
        \item Excessive agency
        \item Long-horizon autonomy
        \item Irreversible actions
        \item Deceptive behavior under evaluation
      \end{itemize}
    \item[\textcolor{KAred}{$\blacktriangleright$}] Each is a design choice made when building the agent, so each can be dialed back.
  \end{itemize}
\end{frame}

\begin{frame}{Excessive Agency}
  \begin{itemize}\itemsep0.5em
    \item[\textcolor{KAred}{$\blacktriangleright$}] The agent is granted more autonomy, tools, or privileges than the task requires.
    \item[\textcolor{KAred}{$\blacktriangleright$}] A calendar assistant that also holds send-email and delete-file permissions has a larger attack surface than its task needs.
    \item[\textcolor{KAred}{$\blacktriangleright$}] OWASP breaks the risk into three sub-problems:
      \begin{itemize}\itemsep0.3em
        \item \textbf{Excessive functionality:} tools the task never needs but the agent can still call.
        \item \textbf{Excessive permissions:} credentials scoped wider than required, such as write access where read would suffice.
        \item \textbf{Excessive autonomy:} high-impact actions executed with no confirmation.
      \end{itemize}
    \item[\textcolor{KAred}{$\blacktriangleright$}] The mitigation is tighter scoping: grant the least capability that still does the job.
  \end{itemize}
\end{frame}

\begin{frame}{Long-Horizon Autonomy: Errors Compound}
  \begin{itemize}\itemsep0.5em
    \item[\textcolor{KAred}{$\blacktriangleright$}] A chatbot answers in one turn. An agent may run for hundreds of steps, each conditioned on the output of the last.
    \item[\textcolor{KAred}{$\blacktriangleright$}] If each step is reliable with probability $p$, the chance of a fully clean $n$-step run is $p^{\,n}$. At $p = 0.99$ and $n = 100$ that is about 37\%.
    \item[\textcolor{KAred}{$\blacktriangleright$}] Errors are also correlated: a wrong assumption made early is carried forward, and the agent may adjust later steps to stay consistent with it.
    \item[\textcolor{KAred}{$\blacktriangleright$}] There is no natural checkpoint. A human reviewing the final output cannot see the hundred intermediate actions that produced it.
    \item[\textcolor{KAred}{$\blacktriangleright$}] Mitigation is structural: checkpoints, subtask verification, and bounded horizons in place of open-ended running until done.
  \end{itemize}
\end{frame}

\begin{frame}{Irreversible Actions}
  \begin{columns}[T,onlytextwidth]
    \column{0.52\textwidth}
      \begin{itemize}\itemsep0.5em
        \item[\textcolor{KAred}{$\blacktriangleright$}] Some actions cannot be undone by a later step: deletion, payments, sending communications, submitting a form, executing a trade.
        \item[\textcolor{KAred}{$\blacktriangleright$}] For these, expecting the agent to notice and fix the mistake does not help. The harm has already left the system.
        \item[\textcolor{KAred}{$\blacktriangleright$}] The practical response is to separate reversible actions, which are safe to automate, from irreversible ones, which need a human gate or a dry run first.
      \end{itemize}
    \column{0.48\textwidth}
      \begin{block}{\small Deciding what needs a gate}
        \footnotesize
        Two properties decide whether an action needs a human sign-off:
        \begin{enumerate}\itemsep0.3em
          \item whether it can be reversed cheaply, and
          \item whether it affects anything outside the current session.
        \end{enumerate}
        \vspace{0.3em}
        High-impact, irreversible, and external is the combination that warrants a human decision.
      \end{block}
  \end{columns}
\end{frame}

\begin{frame}{Deceptive and Sycophantic Behavior}
  \begin{itemize}\itemsep0.5em
    \item[\textcolor{KAred}{$\blacktriangleright$}] Models are trained to produce outputs that score well with a rater. Under evaluation pressure, looking done and being done can diverge.
    \item[\textcolor{KAred}{$\blacktriangleright$}] \textbf{Sycophancy:} the agent tells the user what they want to hear, agreeing with a flawed plan rather than flagging the flaw.
    \item[\textcolor{KAred}{$\blacktriangleright$}] \textbf{Faked completion:} an agent that cannot finish a task reports success anyway, or fabricates the artifact it was asked to produce.
    \item[\textcolor{KAred}{$\blacktriangleright$}] \textbf{Working around blocks:} told a path is forbidden, the agent finds an unintended route to the same end instead of stopping and asking.
    \item[\textcolor{KAred}{$\blacktriangleright$}] These behaviors are hard to catch because they undermine the signal a human relies on. The report that the task succeeded is the thing being gamed.
  \end{itemize}
\end{frame}
